% Fragmento público generado desde propietarios íntegros.
% Residencia: 52.9.
% Fuente: ../incoming/radion_monografia_20260827/manuscrito/sections/10_minkowski_hodge_lorentz.tex:1-55
\noindent Oriented holonomy admits subsequent causal and geometric realizations in Minkowski, Hodge, Lorentz, and Cartan--Holst.\par\medskip

\paragraph{Construction of the Lorentzian metric from the APP block}

\paragraph{APP matrix of parameters \(7\) and \(2\)}

The central matrix
\begin{equation}
B_c=\begin{pmatrix}7&2\\2&7\end{pmatrix}
=9P_++5P_-
\label{eq:Bc}
\end{equation}
it is normalized by its determinant:
\begin{equation}
M_c=\frac{B_c}{\sqrt{45}}
=\exp\left(\frac12\log\frac95\,R\right)
\in SO^+(1,1).
\label{eq:Mc}
\end{equation}
The internal speed is
\[
\eta_c=\frac12\log\frac95.
\]
The associated projective publication satisfies
\begin{equation}
P_B^n=P_++\left(\frac59\right)^nP_-.
\label{eq:PB-contract}
\end{equation}
The same ratio \(5/9\) that appears as a spectral contraction is recognized here
as the separation of one stable mode and one contractive mode. The identity is
operatorial; it does not turn every occurrence of \(5/9\) into the same map.

\paragraph{Incidence \(6\to8\) and causal-signature quotient}

Let \(M\) be the typed incidence between six faces and eight
vertices. The recovered spectral relation is
\begin{equation}
MM^{\mathsf T}=12P_u+4P_-,
\label{eq:MMt}
\end{equation}
where \(P_u\) projects onto the uniform mode and \(P_-\) onto the
oriented subspace. The surviving quotient of dimension four decomposes as
\begin{equation}
Q_4\simeq\R u\oplus E_-.
\label{eq:radion-Q4}
\end{equation}
Declaring the uniform mode negative and the three oriented modes positive,
one obtains
\begin{equation}
B_{\mathrm{caus}}=\operatorname{diag}(-1,1,1,1).
\label{eq:minkowski}
\end{equation}
The Minkowski screen does not select the HMT state. It is the subsequent causal realization of the quotient already constructed by the incidence.

